% ##############################################################################
\section{Conclusion}
\label{sec:conclusion}

We have described \Noesis, a protocol that treats machine intelligence as
a fungible, but not homogeneous, commodity, and that organizes its
exchange around two cooperating components: \NoesisMarket, a periodic
call auction that clears bids and asks bundling capability, latency, and
reliability into a single contract, and \NoesisServer, the gateway that
executes those contracts, meters whether they were fulfilled, and reuses
the resulting traffic for difficulty-aware routing, answer fusion, and
distillation. Each of these two components is, in turn, a thin
coordination layer around independently pluggable sub-components, matching
engine, capability measurement, reputation and feedback, answer fusion, and
pricing dissemination (Section~\ref{sec:modularity}), rather than a fixed
monolith. We formalized the contract and auction notation, stated the
market-clearing problem as a linear program, established existence of a
clearing price under a completeness assumption on the compatibility
graph, and proposed a hypothesis-testing approach to distinguishing
genuine contract violations from sampling noise.

The design suggests that the two problems, price discovery for a bundled
quality good and verified fulfillment of that good, are naturally
complementary: a market without fulfillment monitoring has no way to
enforce what was sold, and a fulfillment layer without a market has no
price signal to feed back into. Section~\ref{sec:decentralized_extension}
sketches one further direction: a blockchain-based, blind-bid variant of
\NoesisMarket{} that trades the operational simplicity of a centralized
clearinghouse for censorship resistance and bidder anonymity, at the cost
of the gas, latency, and stake-sizing questions consolidated in
Section~\ref{sec:open_questions}.

We have not, however, evaluated this
design against real traffic, real providers, or strategic sellers, and
Section~\ref{sec:mechanism_design_risks} makes clear that the mechanism's
robustness to bid shading, capability misrepresentation, and collusion is
an open question rather than a settled property. The open questions of
Section~\ref{sec:open_questions}, including the interface contracts each
pluggable component of Section~\ref{sec:modularity} would need to be
truly swappable, are intended as a concrete, incremental path toward
closing that gap, starting from components, an auction simulator and a
logging proxy, that are each independently useful and testable before the
full closed loop is attempted.
